\chapter*{\chapterstyle{Algebre I {:} Corps des Complexes}}
On note \(\C\) l'extension du corps \(\R\) avec les deux lois usuelles. Une des propriétés fondamentales de \(\C\) est qu'il est algébriquement clos, ie tout polyômes non trivial à coefficient dans \(\C\) admet au moins une racine. \<

Soit \(z, z' \in \C\) et \(\theta \in \R/2\pi\). On notera \(\rho\) le module de \(z\) pour faciliter la lecture et \(a + bi\) la forme algébrique de \(z\).

\section*{\subsecstyle{Module {:}}}
On apelle \textbf{module} de \(z\) le nombre réel positif \(|z|\) tel que :
\[
   |z| = \sqrt{a^2 + b^2} = \sqrt{z\overline{z}} 
\]
C'est un \textbf{prolongement} de la fonction valeur absolue à \(\C\).
On note que c'est une opération \textbf{multiplicative} mais \textbf{pas} additive, en effet on a l'inégalité triangulaire:
\[
   |z + z'| \leq |z| + |z'|   
\]
\section*{\subsecstyle{Conjugé {:}}}
On apelle \textbf{conjugaison} l'automorphisme involutif qui à \(z\) associe son \textbf{conjugué}, noté \(\overline{z}\) tel que:
\[
      \overline{z} := a - bi = \rho\cos\theta - i\sin\theta = \rho e^{-i\theta}
\]
C'est un morphisme \textbf{additif} et \textbf{multiplicatif}.

Graphiquement, le conjugué d'un nombre complexe est \textbf{le symétrique} de \(z\) par rapport à l'axe réel.\<

(\(\ast\)) On a alors la propriété suivante {:}
\[
   \boxed{
      \begin{aligned}
         \Re(z) := \frac{z + \overline{z}}{2}
      \end{aligned}
      \hspace*{40pt}
      \begin{aligned}
         \Im(z) := \frac{z - \overline{z}}{2i}
      \end{aligned}
   }
\]
\section*{\subsecstyle{Forme trigonométrique {:}}}
On montre qu'il existe un angle \(\theta\) tel que:
\[
   z := \rho(\cos\theta+i\sin\theta)   
\]
Cette forme est apellée \textbf{forme trigonométrique} de \(z\) et elle est unique modulo \(2\pi\).\+
Elle se déduit du nombre \(\frac{z}{|z|} \in \U\) et de l'identité trigonométrique fondamentale.\<

On peut noter qu'on apelle \(\theta\) ainsi obtenu \textbf{l'argument} de \(z\).
\section*{\subsecstyle{Forme exponentielle {:}}}
De même pour \(\theta\) l'argument de \(z\), on a:
\[
   z := \rho e^{i\theta}   
\]
Cette forme est apellée \textbf{forme exponentielle} de \(z\) est également unique modulo \(2\pi\).\<

On montre alors les \textbf{formules d'Euler}:
\[
   \boxed{
      \begin{aligned}
         \cos\theta := \frac{e^{i\theta}+e^{-i\theta}}{2}
      \end{aligned}
      \hspace*{40pt}
      \begin{aligned}
         \sin\theta := \frac{e^{i\theta}-e^{-i\theta}}{2i}
      \end{aligned}
   }
\]
Ces formules se déduisent de (\(\ast\)). Elle sont trés importantes car elle permettent de \textbf{linéariser} des expression trigonométriques.\<

Les propriétés de l'exponentielle nous permettent aussi d'obtenir:
\begin{align*}
   \arg(zz') &\eqmod{2\pi} \arg(z) + \arg(z') \\
   \arg(\frac{z}{z'}) &\eqmod{2\pi} \arg(z) - \arg(z')
\end{align*}
\section*{\subsecstyle{Formule de Moivre {:}}}
Un propriété importante des formes trigonométriques et exponentielle est:
\[
   {(e^{i\theta})}^n = e^{n(i\theta)}
\]
De manière analogue, on a:
\[
   {(\cos\theta + i\sin\theta)}^n = \cos n\theta + i\sin n\theta
\]
Ici, on a choisi de considérer \(z \in \U\) mais ces propriétés sont vraies pour \textbf{tout nombre complexe}, il suffit alors de factoriser par le module pour se ramener au cas ci-dessus.
\section*{\subsecstyle{Racines n-ièmes {:}}}
Soit \(n \in \N\), une partie importante des problèmes impliquant des nombres complexes proviennent d'équations d'inconnue \(Z\) de la forme:
\[
   Z^n = z
\]
On peut montrer que l'ensemble des solutions de ce type de problème est:
\[
   \boxed{
      S := \Bigl\{ \sqrt[n]{\rho}e^{i\frac{\theta + 2k\pi}{n}}\; ; \; k \in \bigl\{0, 1, \ldots , n-1 \bigl\} \Bigl\}   
   }
\]
\underline{Cas particulier {:}}
Si on a une racine n-ième \(Z_0\) de \(Z\) et qu'on connait les racines n-ièmes de l'unité, alors on peut obtenir toutes les racines n-ièmes de \(Z\) gràce à:

\[
   \Bigl\{ Z \in \C \; ; \; Z^n = z \Bigl\} = \Bigl\{ Z_0u \; ; \; u \in \U_n \Bigl\}   
\]
\section*{\subsecstyle{Le nombre complexe \(j\) {:}}}
On note \(j\) la première racine troisième de l'unité.
Le nombre \(j\) est singulier, car il vérifie:
\[
   j^2 = j^{-1} = \overline{j}
\]
Graphiquement, on peut observer que les images des nombres \(1\), \(j\) et \(\overline{j}\) forment un triangle équilatéral inscrit dans le cercle trigonométrique.


\chapter*{\chapterstyle{Algebre I {:} Anneau des Polynômes}}
\section*{\subsecstyle{Propriétés générales {:}}}
\section*{\subsecstyle{Degré et Valuation {:}}}
\section*{\subsecstyle{Divisibilité {:}}}
\section*{\subsecstyle{Racines et Factorisation {:}}}

\chapter*{\chapterstyle{Algebre I {:} Espaces Vectoriels}}
\section*{\subsecstyle{Généralités {:}}}
\section*{\subsecstyle{Sous-espace Engendré {:}}}
\section*{\subsecstyle{Intersections et Sommes {:}}}
\section*{\subsecstyle{Familles {:}}}
\section*{\subsecstyle{Bases {:}}}
